%============================================================
% SISTEMATIKA PENULISAN: PROPOSAL PA (Bab 1-3)
% Di-input oleh chapters/bab1.tex via \input{\SistematikaFile}
% Dipanggil dari main_proposal.tex
%============================================================

\vspace{0.3cm}
\begin{description}[leftmargin=2.2cm, labelsep=0cm, itemsep=0.1cm]
  \item[{\makebox[2.2cm][l]{\textbf{Bab 1}}}] \textbf{Pendahuluan}\\
    Bab ini berisi latar belakang yang mendasari dilakukannya proyek akhir, rumusan permasalahan yang dihadapi pada infrastruktur \textit{backend} aplikasi FoodLAB terkait \textit{monitoring} dan \textit{self-healing}, batasan masalah yang membatasi ruang lingkup penelitian, tujuan yang ingin dicapai, manfaat yang diharapkan, dan sistematika penulisan laporan ini.

  \item[{\makebox[2.2cm][l]{\textbf{Bab 2}}}] \textbf{Kajian Pustaka}\\
    Bab ini membahas deskripsi permasalahan pada \textit{backend} FoodLAB, teori penunjang yang menjadi dasar perancangan sistem (meliputi \textit{observability}, \textit{agentic AI}, teknologi NATS, PostgreSQL, Docker, Go, React, dan Keycloak), serta kajian penelitian terkait beserta identifikasi celah penelitian (\textit{research gap}) yang menjadi posisi proyek akhir ini.

  \item[{\makebox[2.2cm][l]{\textbf{Bab 3}}}] \textbf{Desain Sistem}\\
    Bab ini membahas deskripsi solusi, perancangan arsitektur \textit{high-level}, alur \textit{handshake} dan \textit{bootstrap agent}, arsitektur \textit{perception agents}, alur kerja \textit{reasoning engine} (RCA \textit{pipeline}), \textit{clean architecture backend}, alur kerja \textit{code fix agent}, desain antarmuka \textit{dashboard} operator, batasan teknis dan pertimbangan privasi data, serta rancangan skenario pengujian dan evaluasi sistem.
\end{description}
